\section{\texorpdfstring{Full Primitive Higher-\(\psi\) Koszul Homology}{Full Primitive Higher-psi Koszul Homology}}
\label{sec:app-full-psi-homology}

The polynomial PBW one-\(\psi\) descendant labels \(\Psi_{a,b}\) of the
open BV brane sector and the residual Matlis principal parts
\(\rho_{a,b}\) of \S~\ref{sec:pro-matlis-target} share
the bidegree index set \((a,b)\) and carry different
\(\mathfrak h\)-actions.  The defect position is polynomial; its
cotangent momentum is residual.  The separation is the homology of a
single universal calculation: the full primitive cyclic Koszul
homology of \((A_\psi,\,d\psi=xy-yx)\) graded by \(\psi\)-count and by
\((R,S)=(\#x+\#\psi,\#y+\#\psi)\).  The surviving \(H_1\)-line is the
special-fibre image of the PBW Rees lattice; the associated-graded
morphism
\(\operatorname{gr}\Phi:\bar A_{\mathrm{desc}}[1]\to\mc P\) is the
bidegree-only special-fibre relabelling.

The universal primitive single-trace Koszul complex is attached to the
moment-map equation \([x,y]=0\).  Universal means that no
choice of Matlis target, PBW model, finite \(N\) trace identity, or quantum
deformation is used: after stable trace invariant theory has identified
primitive traces with graded cyclic words, the answer is the cellular homology
of a chip-moving state complex.  Every finite-bidegree stable primitive
descendant calculation over \(\C\), with \(d\psi=xy-yx\) and no trace
relations beyond graded cyclicity, is obtained from this one.  Passage to
completed coefficient systems requires an exact coefficient functor, or an
explicit condition excluding new \(\varprojlim^1\) classes.

\begin{defn}[Graded cyclic word complex]
\label{def:app-full-psi-complex}
  Let
	  \[
	    A_\psi=\C\langle x,y,\psi\rangle,\qquad
	    |x|=|y|=0,\quad |\psi|\equiv1\pmod 2,
	  \]
  and let \(d\) be the odd derivation determined by
  \[
    d\psi=xy-yx,\qquad dx=dy=0.
  \]
  Put
  \[
    \operatorname{Cyc}(A_\psi)
      =A_\psi/[A_\psi,A_\psi]_{\mathrm{gr}},
      \qquad [uv]=(-1)^{|u||v|}[vu].
  \]
  Give the letters weights
  \[
    \operatorname{wt}(x)=(1,0),\qquad
    \operatorname{wt}(y)=(0,1),\qquad
    \operatorname{wt}(\psi)=(1,1).
  \]
  For \(R,S\geq0\) define
  \[
    K^p_{R,S}
      =
      \C\{[w]_{\mathrm{cyc}}:
        \#\psi(w)=p,\ \#x(w)=R-p,\ \#y(w)=S-p\},
  \]
  with \(K^p_{R,S}=0\) unless \(0\leq p\leq\min(R,S)\).  If
  \(w=a_1\cdots a_n\), the boundary is
  \[
  \begin{aligned}
    \partial[w]
      =
      \sum_{i:a_i=\psi}(-1)^{\#\{j<i:a_j=\psi\}}
      [&a_1\cdots a_{i-1}xy\,a_{i+1}\cdots a_n\\
       &-a_1\cdots a_{i-1}yx\,a_{i+1}\cdots a_n]_{\mathrm{cyc}} .
  \end{aligned}
  \]
  Thus \(\partial:K^p_{R,S}\to K^{p-1}_{R,S}\).
\end{defn}

\begin{lem}[Stable super-cyclic trace presentation]
\label{lem:app-full-stable-super-traces}
  Fix a total word length \(d\).  In the stable range \(N\geq d\), the
  primitive \(GL_N\)-invariant single-trace functions in the variables
  \(x,y,\psi\), homogeneous of total word length \(d\), are exactly the
  graded cyclic words in \(\operatorname{Cyc}(A_\psi)\).  Equivalently, after
  connected single-trace extraction, the only stable relation among words in
  \(x,y,\psi\) is
  \[
    \operatorname{Tr}(uv)=(-1)^{|u||v|}\operatorname{Tr}(vu).
  \]
\end{lem}

\begin{proof}
  First replace the odd letter \(\psi\) by \(p\) distinct even marked
  matrices \(\eta_1,\ldots,\eta_p\), and take the summand multilinear in the
  \(\eta_i\)'s and homogeneous of the prescribed \(x,y\)-degree.  The
  Procesi--Razmyslov stable invariant theorem
  \cites{procesi,razmyslov} identifies the primitive single-trace invariants
  of total word length \(d\), for \(N\geq d\), with ordinary cyclic words in
  the letters \(x,y,\eta_1,\ldots,\eta_p\), modulo only cyclicity of the
  trace.

  The symmetric group \(\mathbb S_p\) permutes the marked letters.  Passing
  from \(p\) even marked letters to one odd letter is the sign-isotypic
  projection for this action, followed by the identification
  \(\eta_i\mapsto\psi\).  Since the ground field has characteristic zero, this
  projection is exact.  Interchanging two marked occurrences contributes the
  sign of the transposition, and rotating a block with \(a\) marked letters
  past a block with \(b\) marked letters contributes \((-1)^{ab}\).  This is
  precisely the graded cyclic rule for a single odd letter \(\psi\).  No other
  stable trace relation remains in word length \(d\).
\end{proof}

\begin{lem}[Boundary signs and \(\partial^2=0\)]
\label{lem:app-full-boundary-d2}
  The formula of Definition~\ref{def:app-full-psi-complex} is well-defined on
  the graded cyclic quotient and satisfies \(\partial^2=0\).
\end{lem}

\begin{proof}
  On the free algebra \(A_\psi\), the displayed formula is the derivation
  rule for the odd derivation \(d\): replacing the \(i\)-th occurrence of
  \(\psi\) crosses exactly the preceding odd letters and hence contributes
  \((-1)^{\#\{j<i:a_j=\psi\}}\).  Since \(d(x)=d(y)=0\) and
  \(d(\psi)=xy-yx\) is even and \(d\)-closed, one has \(d^2=0\) on the
  generators, hence on all of \(A_\psi\).

  If \(u,v\) are homogeneous, then
  \[
    d(uv-(-1)^{|u||v|}vu)
  \]
  is a sum of graded commutators, using the odd Leibniz rule and the fact that
  \(d\) lowers the number of odd letters by one.  Thus \(d\) preserves
  \([A_\psi,A_\psi]_{\mathrm{gr}}\), so it descends to
  \(\operatorname{Cyc}(A_\psi)\).  The descended operator is exactly
  \(\partial\), and \(\partial^2=0\) follows from \(d^2=0\).
\end{proof}

\begin{lem}[Exact \(p\)-face formula]
\label{lem:app-full-p-face-formula}
  Let \(u_1,\ldots,u_p\) be ordinary words in \(x,y\).  In
  \(K^{p-1}_{R,S}\),
  \[
  \begin{aligned}
    \partial[\psi u_1\psi u_2\cdots\psi u_p]
      =
      \sum_{j=1}^p(-1)^{j-1}
      [&\psi u_1\cdots\psi u_{j-1}xy\,u_j\psi u_{j+1}\cdots\psi u_p\\
       &-\psi u_1\cdots\psi u_{j-1}yx\,u_j\psi u_{j+1}\cdots\psi u_p]_{\mathrm{cyc}},
  \end{aligned}
  \]
  where, for \(j=1\), the initial displayed \(\psi u_1\cdots\psi u_{j-1}\)
  is absent.
\end{lem}

\begin{proof}
  This is the boundary of Definition~\ref{def:app-full-psi-complex} applied
  to the representative \(\psi u_1\psi u_2\cdots\psi u_p\).  The \(j\)-th
  displayed \(\psi\) has \(j-1\) preceding odd letters, and replacing it by
  the even word \(xy-yx\) produces no further Koszul sign.
\end{proof}

\begin{lem}[Two-\(\psi\) boundary sign]
\label{lem:app-full-two-psi-boundary-sign}
  For ordinary words \(u,v\),
  \[
    \partial[\psi u\psi v]
      =[\psi vxyu]-[\psi vyxu]-[\psi uxyv]+[\psi uyxv].
  \]
  After removing the common front \(\psi\) from each term, this is the
  boundary
  \[
    vxyu-vyxu-uxyv+uyxv.
  \]
\end{lem}

\begin{proof}
  Lemma~\ref{lem:app-full-p-face-formula} with \(p=2\) gives
  \[
    [(xy-yx)u\psi v]-[\psi u(xy-yx)v].
  \]
  The first two terms are cyclically rotated to put the remaining \(\psi\) in
  front:
  \[
    [(xy)u\psi v]=[\psi vxyu],\qquad
    [(yx)u\psi v]=[\psi vyxu].
  \]
  The moved block is even, so no Koszul sign is produced.  The second
  replacement already has the surviving \(\psi\) in front, and its overall
  sign is the sign \(-1\) from the preceding odd letter.  Expanding gives the
  displayed formula.

  Interchanging \(u\) and \(v\) gives
  \(uxyv-uyxv-vxyu+vyxu\), the negative of the displayed expression.  Thus
  the formula is compatible with the alternating two-\(\psi\) quotient used
  below.
\end{proof}

\begin{constr}[One-\(\psi\) three-term skeleton]
\label{constr:app-one-psi-three-term-skeleton}
  For \(k,l\geq0\), let \(\mathsf W_{k,l}\) be the ordinary words with
  \(k\) letters \(x\) and \(l\) letters \(y\), and let
  \(\mathsf N_{a,b}\) be the necklaces with \(a\) letters \(x\) and
  \(b\) letters \(y\).  The one-\(\psi\) skeleton in bidegree \((k,l)\) is
  the finite complex
  \[
    C_2(k,l)\xrightarrow{\partial_2}C_1(k,l)
      \xrightarrow{\partial_1}C_0(k,l),
  \]
  with
  \[
    C_1(k,l)=\C\{\mathsf W_{k,l}\},\qquad
    C_0(k,l)=\C\{\mathsf N_{k+1,l+1}\},
  \]
  and
  \[
    C_2(k,l)=
    \bigl(\C\{(u,v):\deg u+\deg v=(k-1,l-1)\}\bigr)_{\mathrm{alt}},
  \]
  where \((u,v)=-(v,u)\) and \((u,u)=0\).  If \(k=0\) or \(l=0\), then
  \(C_2(k,l)=0\).  The boundaries are
  \[
    \partial_1(w)=[xyw]_{\mathrm{cyc}}-[yxw]_{\mathrm{cyc}},
  \]
  and
  \[
    \partial_2(u,v)=vxyu-vyxu-uxyv+uyxv.
  \]
  This is the truncation
  \(K^2_{k+1,l+1}\to K^1_{k+1,l+1}\to K^0_{k+1,l+1}\) after putting the
  remaining \(\psi\) at the front of a one-\(\psi\) cyclic word.
  Equivalently, it is the cellular chain complex in dimensions
  \(2,1,0\) of the state complex constructed below: \(C_0\) records chip
  configurations, \(C_1\) records one moving chip, and \(C_2\) records two
  independent moving chips with the orientation fixed by the ordered moving
  edges.
\end{constr}

\begin{proof}
  The identifications of \(C_1\) and \(C_0\) are exactly the graded cyclic
  normal forms in \(\psi w\) and the ordinary necklaces.  The source
  \(C_2\) is the two-\(\psi\) summand of \(K^2_{k+1,l+1}\): after cutting
  at one displayed \(\psi\), a class is written \(\psi u\psi v\), and
  moving the second \(\psi\)-block past the first contributes the sign
  \(-1\).  Lemma~\ref{lem:app-full-two-psi-boundary-sign} gives
  \(\partial_2\), while the definition of \(\partial\) gives
  \(\partial_1\).
\end{proof}

\begin{constr}[Labelled chip-moving state complex]
\label{constr:app-chip-state-complex}
  Assume \(R,S>0\).  Let \(\mathscr C_S\) be the oriented circle subdivided
  into \(S\) labelled oriented edges and \(S\) labelled vertices.  Define
  \(\widetilde Y_{R,S}\) as follows.  A \(p\)-cell is a state with
  \(R-p\) stationary indistinguishable chips on vertices and \(p\) moving
  chips in the interiors of \(p\) distinct edges.  Several stationary chips
  may occupy the same vertex.  No edge contains more than one moving chip.

  Encode such a state by traversing \(\mathscr C_S\).  At a vertex write one
  \(x\) for each stationary chip sitting there.  Then write \(y\) if the
  following edge is empty and \(\psi\) if the following edge contains a
  moving chip.  This gives a word with \(R-p\) letters \(x\), \(S-p\) letters
  \(y\), and \(p\) letters \(\psi\), with a distinguished starting edge.

  The \(p\)-cell is oriented by the increasing order of the labelled edges
  carrying moving chips.  Its \(j\)-th oriented coordinate has boundary
  \[
    \text{terminal face }xy-\text{ initial face }yx.
  \]
  Hence the cellular boundary in the labelled complex is
  \[
    \sum_j(-1)^{j-1}(\text{replace the \(j\)-th moving edge by \(xy\)}
      -\text{replace it by \(yx\)}).
  \]
  The cyclic group \(\Gamma_S=\Z/S\Z\) rotates the edge labels and acts
  cellularly on \(\widetilde Y_{R,S}\).
\end{constr}

\begin{lem}[Cellular chain identification with cyclic relabelling]
\label{lem:app-state-chain-identification}
  For \(R,S>0\), the cellular chain complex of
  \(\widetilde Y_{R,S}\), after taking \(\Gamma_S\)-coinvariants, is naturally
  isomorphic to \(K^\bullet_{R,S}\):
  \[
    C_p(\widetilde Y_{R,S};\C)_{\Gamma_S}\cong K^p_{R,S}.
  \]
  Under this isomorphism the cellular boundary is the boundary
  \(\partial\) of Definition~\ref{def:app-full-psi-complex}.
\end{lem}

\begin{proof}
  A labelled \(p\)-cell is equivalently a labelled cyclic sequence
  \[
    x^{a_1}m_1x^{a_2}m_2\cdots x^{a_S}m_S,\qquad
    m_i\in\{y,\psi\},
  \]
  with exactly \(p\) of the \(m_i\)'s equal to \(\psi\) and
  \(\sum_i a_i=R-p\).  Every graded cyclic word in
  \(K^p_{R,S}\) has a representative of this form: move the cut, if
  necessary, past even \(x\)-letters until it lies immediately before an edge
  marker \(y\) or \(\psi\).  This representative is unique after a starting
  edge has been chosen.  The remaining ambiguity is rotation of the \(S\) edge
  markers, hence quotienting by \(\Gamma_S\) is exactly passage to graded
  cyclic words.  No further relation is introduced: in the labelled complex
  the \(S\) edge markers are ordered, and in the coinvariant complex the only
  identifications are their rotations with the orientation character below.

  The sign is determined as follows.  Suppose a rotation moves a block containing
  \(a\) occurrences of
  \(\psi\) past the complementary block containing \(b\) occurrences of
  \(\psi\).  On cellular orientations this rotation moves \(a\) oriented
  interval coordinates past \(b\) oriented interval coordinates, hence changes
  the orientation by \((-1)^{ab}\).  This is exactly the graded cyclic sign
  \([uv]=(-1)^{|u||v|}[vu]\), since \(x\) and \(y\) are even and every
  \(\psi\) is odd.

  With this orientation convention, the cellular boundary in
  Construction~\ref{constr:app-chip-state-complex} is
  \[
    \sum_j(-1)^{j-1}([xy]-[yx])
  \]
  in the \(j\)-th moving edge.  The terminal face is the state in which the
  chip has crossed the edge and is read as \(xy\); the initial face is read as
  \(yx\).  This is exactly Lemma~\ref{lem:app-full-p-face-formula}, term by
  term.
\end{proof}

\begin{lem}[Abel map for labelled states]
\label{lem:app-abel-map-state-homology}
  For \(R,S>0\), the labelled state complex \(\widetilde Y_{R,S}\) has the
  homology of a circle:
  \[
    H_q(\widetilde Y_{R,S};\C)\cong
    \begin{cases}
      \C,&q=0,\\
      \C,&q=1,\\
      0,&q\geq2.
    \end{cases}
  \]
  The cyclic rotation action of \(\Gamma_S\) is trivial on these two homology
  lines.
\end{lem}

\begin{proof}
  Work in the universal cover of \(\mathscr C_S\), whose vertices are the integers and
  whose edges are the unit intervals.  A lifted labelled cell is described by
  the positions of the \(R\) chips: stationary chips have integer positions,
  and a moving chip on the edge \([m,m+1]\) has coordinate \(m+t\) with
  \(0<t<1\).  The closure of the cell permits \(t=0,1\), which are exactly the
  two vertex states \(yx\) and \(xy\) in the boundary convention of
  Construction~\ref{constr:app-chip-state-complex}.  Define the Abel map
  \[
    A:\widetilde Y_{R,S}\longrightarrow \R/S\Z
  \]
  by summing the chip positions on the oriented circle, counted with
  multiplicity, and reducing modulo \(S\).  On every cell this map is affine.

  Lift the target to \(\R\) and fix a cyclic order of the \(R\) lifted chips.
  A lifted fibre is described by weakly ordered coordinates
  \[
    t_1\leq t_2\leq\cdots\leq t_R\leq t_1+S
  \]
  with fixed sum.  This is a convex polytope.  Its faces are cut out by the
  equalities \(t_i\in\Z\) and by the choices of those unit intervals which
  contain a moving chip; these faces are precisely the cells of the
  chip-moving subdivision meeting the fibre.  Hence every lifted fibre is a
  contractible finite polyhedral complex.  The map is a proper PL map of
  finite polyhedra with acyclic fibres.  The Vietoris--Begle theorem, applied
  to this affine cellular map, gives a homology equivalence
  \[
    H_\bullet(\widetilde Y_{R,S};\C)\cong H_\bullet(S^1;\C).
  \]

  A cyclic relabelling of the \(S\) edges adds a constant to every chip
  coordinate, so it rotates the target circle by an orientation-preserving
  rotation.  Therefore \(\Gamma_S\) acts trivially on
  \(H_0(S^1;\C)\) and \(H_1(S^1;\C)\), and hence on the two homology lines of
  \(\widetilde Y_{R,S}\).
\end{proof}

\begin{lem}[Cyclic Abel quotient and stabilizer fibres]
\label{lem:app-cyclic-abel-stabilizers}
  Let \(Y_{R,S}\) denote the cyclic relabelling quotient of
  \(\widetilde Y_{R,S}\), with the orientation character encoded by the
  graded cyclic rule.  Then the quotient chain complex
  \(C_\bullet(\widetilde Y_{R,S};\C)_{\Gamma_S}\) has the homology of \(S^1\).
  Moreover, every stabilizer quotient of an Abel fibre is contractible.
\end{lem}

\begin{proof}
  By Lemma~\ref{lem:app-abel-map-state-homology}, the labelled Abel map
  is a homology equivalence from \(\widetilde Y_{R,S}\) to the circle.  A
  cyclic relabelling by \(r\in \Gamma_S\) adds \(r\) to each of the \(R\) chip
  positions, hence translates the Abel coordinate by \(Rr\) in
  \(\R/S\Z\).  The action on the target circle is by
  orientation-preserving rotations, so it is trivial on the two homology
  lines.  Since \(\C[\Gamma_S]\) is semisimple,
  \[
    H_\bullet(C_\bullet(\widetilde Y_{R,S};\C)_{\Gamma_S})
      \cong H_\bullet(\widetilde Y_{R,S};\C)_{\Gamma_S}
      \cong H_\bullet(S^1;\C).
  \]

  Now fix an Abel value and a stabilizer subgroup \(G\subset \Gamma_S\) of that
  value.  In the universal cover, choose the lifted fibre before decomposing
  it into chip cells.  It is a convex polytope cut out by weak ordering, the
  fixed-sum equation, and the chip-cell inequalities, and \(G\) acts affinely
  on it.  Averaging a point over \(G\) gives a \(G\)-fixed point.  Straight-line
  contraction to this fixed point is \(G\)-equivariant and respects the
  chip-cell subdivision.  Hence the stabilizer quotient of the Abel fibre is
  contractible.
\end{proof}

\begin{lem}[\(S^1\) generator]
\label{lem:app-full-s1-generator}
  For \(R,S>0\), the cycle
  \[
    \Psi_{R-1,S-1}
      =
      \sum_{w\in\mathsf W_{R-1,S-1}}[\psi w]_{\mathrm{cyc}}
  \]
  represents the fundamental \(H_1\)-line of the cyclic chip-moving
  complex.
\end{lem}

\begin{proof}
  Its boundary is
  \[
    \partial\Psi_{R-1,S-1}
      =
      \sum_{w\in\mathsf W_{R-1,S-1}}
      ([xyw]_{\mathrm{cyc}}-[yxw]_{\mathrm{cyc}}).
  \]
  For a fixed necklace with \(R\) letters \(x\) and \(S\) letters \(y\),
  the coefficient is the number of cyclic \(xy\)-transitions minus the
  number of cyclic \(yx\)-transitions.  These two numbers are equal on a
  circle; hence \(\Psi_{R-1,S-1}\) is closed.

  Let \(\epsilon:K^1_{R,S}\to\C\) send every one-\(\psi\) cyclic word to
  \(1\).  Lemma~\ref{lem:app-full-two-psi-boundary-sign} shows that every
  two-\(\psi\) boundary has two positive and two negative terms, so it
  lies in \(\ker\epsilon\).  But
  \[
    \epsilon(\Psi_{R-1,S-1})=\binom{R+S-2}{R-1}\neq0.
  \]
  By Lemma~\ref{lem:app-cyclic-abel-stabilizers}, the \(H_1\)-line is
  one-dimensional.  Moreover, every one-cell in the displayed sum is oriented
  in the positive Abel direction: its terminal face has the moved chip one edge
  farther along the oriented circle.  Thus the displayed cycle maps to
  \(\binom{R+S-2}{R-1}\) times the oriented circle class and spans the
  fundamental \(H_1\)-line.
\end{proof}

\begin{thm}[Full primitive higher-\(\psi\) Koszul homology]
\label{thm:app-full-primitive-psi-homology}
  For every \(R,S\geq0\),
  \[
    H_p(K_{R,S})\cong
    \begin{cases}
      \C\cdot[x^Ry^S],&p=0,\\
      \C\cdot[\Psi_{R-1,S-1}],&p=1,\ R,S\geq1,\\
      0,&p\geq2,
    \end{cases}
  \]
  where
  \[
    \Psi_{R-1,S-1}
      =
      \sum_{w\in\mathsf W_{R-1,S-1}}[\psi w]_{\mathrm{cyc}} .
  \]
  Consequently the primitive stable Koszul homology is
  \[
    \C[x,y]\oplus\C[x,y]\theta,\qquad
    |\theta|=1,\quad \operatorname{wt}(\theta)=(1,1),
  \]
  with the scalar Hamiltonian removed later if one passes to the reduced
  Hamiltonian sector.
\end{thm}

\begin{proof}
  If \(R=0\) or \(S=0\), then \(K^p_{R,S}=0\) for \(p>0\), and
  \(K^0_{R,S}\) is spanned by the unique cyclic word \(x^R\) or \(y^S\)
  (including the empty word when \(R=S=0\)).  This gives the stated homology.

  Assume \(R,S>0\).  By Lemma~\ref{lem:app-state-chain-identification},
  \(K^\bullet_{R,S}\) is the \(\Gamma_S\)-coinvariant cellular chain complex of
  \(\widetilde Y_{R,S}\).  Since \(\C[\Gamma_S]\) is semisimple, taking
  coinvariants is an exact functor on finite-dimensional \(\C[\Gamma_S]\)-modules.
  Therefore
  \[
    H_p(K_{R,S})
      \cong H_p(\widetilde Y_{R,S};\C)_{\Gamma_S}.
  \]
  Lemma~\ref{lem:app-cyclic-abel-stabilizers} identifies this homology
  with \(H_\bullet(S^1;\C)\).  Thus \(H_p(K_{R,S})\) is one-dimensional
  for \(p=0,1\) and vanishes for \(p\geq2\).
  In particular, the three-term skeleton
  \(C_2\to C_1\to C_0\) already computes \(H_1\): cells of dimension
  \(p\ge3\) have boundaries in \(C_{p-1}\) and cannot change
  \(\ker(C_1\to C_0)/\operatorname{im}(C_2\to C_1)\).

  The degree-zero line is generated by any necklace of multidegree \((R,S)\).
  The functional sending every zero-\(\psi\) necklace to \(1\) kills every
  one-\(\psi\) boundary, so this line is nonzero; it is denoted
  \(\C\cdot[x^Ry^S]\).

  Lemma~\ref{lem:app-full-s1-generator} identifies the degree-one
  generator as \(\Psi_{R-1,S-1}\).
\end{proof}

\begin{rmk}[Numerical anchor]
\label{rmk:app-full-psi-numerical-anchor}
  The script
  \texttt{scripts/check\_one\_psi\_homology.py} tabulates the same
  three-term skeleton over \(\Q\), confirms
  \(\partial_1\partial_2=0\), checks that \(\Psi_{R-1,S-1}\) is closed
  and independent modulo boundaries, and verifies
  \[
    \dim\,H_1(K_{R,S})=1
  \]
  in each of the \(36\) bidegrees \(0\leq R-1,S-1\leq5\).  The same
  script verifies the open Hamiltonian descendant action of
  Proposition~\ref{prop:open-descendant-action} in \(240\) cases with
  exponents at most \(3\), and the closed coadjoint Taylor-dual / residue
  formula of Proposition~\ref{prop:app-matlis-coadjoint-action} in
  \(1225\) cases with exponents at most \(5\).  The
  PBW-versus-Moyal sweep records \(1200\) explicit mismatches and
  \(958\) cases where the coadjoint Moyal action carries genuine
  \(\hbar^{\geq2}\) higher-order terms; the \(70\) linear-generator
  leading-order tests with labels at most \(5\) pin the exact obstruction
  to any \(\hbar\)-adic deformation of the classical
  \(\Psi\to\rho\) projection.  The displayed homology lines and the
  separation in
  Theorem~\ref{thm:pbw-vs-deformation} are therefore proved at the
  bidegree level by an independent finite computation, and not by a
  symbolic shortcut.
\end{rmk}

\begin{constr}[Special-fibre PBW one-\(\psi\) label basis]
\label{constr:app-full-pbw-special-fibre-basis}
  Let
  \[
    \mathfrak h=\C[[z_1,z_2]]/\C\cdot1,\qquad
    \mathfrak h^\vee_{\mathrm{cont}}\cong \mc P
      =\bigoplus_{a+b>0}\C\,\rho_{a,b},
  \]
  with the Hamiltonian Lie algebra of
  Definition~\ref{def:app-matlis-topology} and the Matlis principal-part
  module of Definition~\ref{def:app-matlis-principal-parts}.  Let
  \[
    \mathfrak g
      =\mathfrak h\ltimes\mathfrak h^\vee_{\mathrm{cont}}[1]
  \]
  be the semidirect Hamiltonian shifted-cotangent Lie algebra of
  Proposition~\ref{prop:hamiltonian-polyvector-reduction}, and let
  \[
    U_\hbar(\mathfrak g)
      \supset
    \mathbf S_\hbar(\mathfrak h\oplus\mathfrak h^\vee[1])
  \]
  denote its Rees PBW lattice over \(\C[[\hbar]]\), namely the Rees
  algebra of the canonical PBW filtration.  The special fibre is
  \[
    \mathbf S_\hbar(\mathfrak h\oplus\mathfrak h^\vee[1])
      \big|_{\hbar=0}
      \cong
    \mathbf S(\mathfrak h\oplus\mathfrak h^\vee[1]),
  \]
  the symmetric algebra carrying the trivial \(\mathfrak h\)-action on
  generators.  The \emph{primitive one-\(\psi\) PBW label sector} is the
  one-tensor-factor piece in \(\mathfrak h^\vee[1]\), tensored with the
  symmetric algebra of \(\mathfrak h\), in bidegree \((a,b)\):
  \[
    \bar A_{\mathrm{desc}}[1]
      :=
    \bigoplus_{a+b>0}\C\cdot e^{\mathrm{PBW}}_{a,b}\,[1],
    \qquad
    e^{\mathrm{PBW}}_{a,b}\,[1]\in
    \mathbf S^1(\mathfrak h^\vee[1])\otimes
    \mathbf S(\mathfrak h)\big|_{\hbar=0} .
  \]
  The boundary pairing of
  Corollary~\ref{cor:cotangent-boundary-pairing} identifies this label
  sector with the surviving primitive
  \(H_1\)-lines of
  Theorem~\ref{thm:app-full-primitive-psi-homology}: under
  \[
    e^{\mathrm{PBW}}_{a,b}\,[1]
      \longmapsto
    [\Psi_{a,b}],
    \qquad a+b>0,
  \]
  the \(\Psi_{a,b}\) of Theorem~\ref{thm:app-full-primitive-psi-homology}
  are exactly the special-fibre representatives of the PBW label basis,
  with the binomial weight \(\epsilon(\Psi_{a,b})=\binom{a+b}{a}\)
  recording the totally symmetric averaging of the cyclic words.
\end{constr}

\begin{prop}[Associated-graded label morphism]
\label{prop:app-full-gr-Phi}
  Let
  \[
    \widetilde\Psi_{a,b}
      :=
    \binom{a+b}{a}^{-1}\Psi_{a,b}
      \in\bar A_{\mathrm{desc}}[1]
  \]
  be the unit-normalised special-fibre PBW label, so that
  \(\epsilon(\widetilde\Psi_{a,b})=1\), and let
  \(\rho_{a,b}\in\mc P\) be the Matlis principal-part basis of
  Definition~\ref{def:app-matlis-principal-parts}.  The bidegree-only
  relabelling
  \[
    \operatorname{gr}\Phi:
    \bar A_{\mathrm{desc}}[1]
      \xrightarrow{\;\;\sim\;\;}
    \mc P,
    \qquad
    \widetilde\Psi_{a,b}\longmapsto\rho_{a,b},\quad a+b>0,
  \]
  is an isomorphism of \(\Z_{\geq0}^2\)-graded
  \(\C\)-vector spaces.  It is the special-fibre associated-graded
  label match of
  Theorem~\ref{thm:app-matlis-rees-fourier-bridge} in the
  \(\Psi\)-presentation of
  Construction~\ref{constr:app-full-pbw-special-fibre-basis}.

  This morphism is \emph{not} an \(\mathfrak h\)-module isomorphism.
  More precisely, on the linear Hamiltonians,
  \[
    z_1\cdot\widetilde\Psi_{a,b}
      =b\,\widetilde\Psi_{a,b-1},
      \qquad
    z_2\cdot\widetilde\Psi_{a,b}
      =-a\,\widetilde\Psi_{a-1,b},
  \]
  while
  \[
    z_1\cdot\rho_{a,b}
      =-(b+1)\,\rho_{a,b+1},
      \qquad
    z_2\cdot\rho_{a,b}
      =(a+1)\,\rho_{a+1,b}.
  \]
  The two actions have opposite shift directions and unequal scalars; in
  particular the polynomial PBW action is locally nilpotent on every
  \(\widetilde\Psi_{a,b}\), whereas the principal-part coadjoint action
  is not locally nilpotent on any nonzero \(\rho_{a,b}\).  The
  obstruction to an \(\hbar\)-adic same-action lift of
  \(\operatorname{gr}\Phi\) is the local-nilpotence dichotomy of
  Theorem~\ref{thm:app-matlis-polynomial-descendant-obstruction} and
  Corollary~\ref{cor:app-matlis-no-flat-same-action-interpolation}.
\end{prop}

\begin{proof}
  The bidegree component of \(\bar A_{\mathrm{desc}}[1]\) is the line
  \(\C\cdot e^{\mathrm{PBW}}_{a,b}[1]\) for \(a+b>0\); the bidegree
  component of \(\mc P\) is \(\C\cdot\rho_{a,b}\), \(a+b>0\), by
  Definition~\ref{def:app-matlis-principal-parts}.  The displayed map
  is the bidegree identity, hence an isomorphism of bigraded vector
  spaces.

  The polynomial Hamiltonian action on \(\widetilde\Psi_{a,b}\) is
  Proposition~\ref{prop:open-descendant-action}: for
  \(F_{p,q}=\operatorname{Tr}(\phi_1^p\phi_2^q)\),
  \(F_{p,q}\cdot\widetilde\Psi_{a,b}=(pb-qa)\widetilde\Psi_{a+p-1,b+q-1}\).
  Setting \((p,q)=(1,0),(0,1)\) gives the displayed two formulas.  The
  principal-part coadjoint action is
  Proposition~\ref{prop:app-matlis-coadjoint-action}: for
  \((p,q)=(1,0),(0,1)\) the formula
  \(z_1^pz_2^q\cdot\rho_{a,b}=-(pb-qa+p-q)\rho_{a-p+1,b-q+1}\) gives the
  two displayed formulas with the principal-part shifts
  \(b+1\) and \(a+1\).

  For local nilpotence: \(z_1\) lowers the second index of
  \(\widetilde\Psi_{a,b}\) by one and is killed by \(\widetilde\Psi_{a,0}\)
  for \(a\geq0\); a finite linear combination is exhausted after at most
  \(\max b+1\) iterations.  By contrast, \(z_1\) raises the second index
  of \(\rho_{a,b}\) and produces a nonzero rising-factorial coefficient
  \((-1)^N(b+1)(b+2)\cdots(b+N)\rho_{a,b+N}\), which is never zero.  The
  same dichotomy holds for \(z_2\).  Local nilpotence is an isomorphism
  invariant; therefore no module isomorphism with the displayed
  \(\mathfrak h\)-actions can exist.

  The obstruction to extending \(\operatorname{gr}\Phi\) to an
  \(\hbar\)-adic same-action map is exactly this local-nilpotence
  dichotomy, and is recorded as
  Theorem~\ref{thm:app-matlis-polynomial-descendant-obstruction} together
  with the deformation-level
  Corollary~\ref{cor:app-matlis-no-flat-same-action-interpolation}.
  Theorem~\ref{thm:app-matlis-rees-fourier-bridge} supplies the
  positive Fourier--Rees bridge that does identify the two label sets,
  but only after the Fourier twist of the Hamiltonian action.
\end{proof}

\begin{rmk}[Smallest explicit mismatch: \((p,q,a,b)=(4,0,1,1)\)]
\label{rmk:app-full-instant-computation}
  The polynomial PBW action on the descendant labels and the Moyal
  coadjoint action on the principal parts disagree already at one
  exponent pair recorded by the verification script.  Take the quartic
  Hamiltonian \(F_{4,0}=\operatorname{Tr}(\phi_1^4)\) and the
  bidegree-\((1,1)\) descendant target.  On the polynomial PBW side
  Proposition~\ref{prop:open-descendant-action} gives
  \[
    F_{4,0}:\widetilde\Psi_{1,1}\longmapsto
      (4\cdot1-0\cdot1)\,\widetilde\Psi_{1+4-1,1-1}
      =4\,\widetilde\Psi_{4,0}.
  \]
  Direct calculation in the Moyal coadjoint module gives
  \[
    \operatorname{ad}^*_{z_1^4}(\delta_{1,1})
      =-24\,\hbar^2\,\delta_{0,4},
  \]
  the only nonzero term of order \(\leq3\) in the formal Moyal expansion.
  No Fourier-untwisted relabelling of the bidegree pair \((1,1)\to(0,4)\)
  versus \((1,1)\to(4,0)\) carries the polynomial PBW action onto the
  Moyal coadjoint action.  Replacing the Moyal coadjoint by its classical
  truncation \((\hbar^0)\) gives
  \(\operatorname{ad}^*_{z_1^4}(\rho_{1,1})=0\) by the principal-part
  formula at this exponent, again disagreeing with the polynomial PBW
  output \(4\widetilde\Psi_{4,0}\).  The script confirms this fingerprint
  exactly and embeds it inside a sweep of \(1200\) such mismatches.
\end{rmk}

\begin{thm}[Universal primitive descendant calculation]
\label{thm:app-full-universal-primitive-descendant}
  Let \(\mathsf K_{\mathrm{univ}}\) be the stable primitive cyclic Koszul
  complex generated by two even letters \(x,y\) and one odd letter
  \(\psi\), with \(d\psi=xy-yx\), decomposed by the weights of
  Definition~\ref{def:app-full-psi-complex}.  Then
  \[
    H_\bullet(\mathsf K_{\mathrm{univ}})
      \cong
      \C[x,y]\oplus\C[x,y]\theta,\qquad
      |\theta|=1,\quad\operatorname{wt}(\theta)=(1,1),
  \]
  with \(x^ay^b\theta\) represented by the symmetrised primitive cycle
  \(\Psi_{a,b}\).  This is the universal primitive calculation in the
  following precise sense.  If \(E_\bullet\) is any stable primitive
  single-trace model whose generators are the images of \(x,y,\psi\), whose
  differential sends \(\psi\) to \(xy-yx\), and whose only stable primitive
  trace relation is graded cyclicity, then the induced map
  \[
    \mathsf K_{\mathrm{univ}}\longrightarrow E_\bullet
  \]
  factors on homology through the displayed two-line module.  Hence there
  are no hidden universal primitive higher-\(\psi\) classes from which a
  further cotangent or descendant bridge can be built.
\end{thm}

\begin{proof}
  The first assertion is Theorem~\ref{thm:app-full-primitive-psi-homology}
  summed over all bidegrees \((R,S)\), with
  \(x^ay^b\theta=\Psi_{a,b}\) and \(R=a+1\), \(S=b+1\) in degree one.
  For the universal property,
  Lemma~\ref{lem:app-full-stable-super-traces} says that the stable
  primitive trace span is exactly the graded cyclic quotient of the free
  algebra on \(x,y,\psi\).  Therefore any stable primitive model satisfying
  the stated hypotheses receives a unique chain map from
  \(\mathsf K_{\mathrm{univ}}\) after choosing the images of \(x,y,\psi\).
  The homology of the source is the two-line module just computed.  Since
  \(H_p(K_{R,S})=0\) for \(p\ge2\) in every bidegree, every universal
  primitive class in the target comes from the degree-zero polynomial line
  or the degree-one \(\Psi\)-line.  The statement is independent of the
  later choice of principal-part, PBW, or Moyal coefficient module.
\end{proof}

\begin{thm}[Hamiltonian-sector extraction]
\label{thm:app-full-hamiltonian-sector-extraction}
  The full primitive homology decomposes as
  \[
    \C\cdot[1]\oplus \bar A
      \oplus \C\cdot[\Psi_{0,0}]
      \oplus \bar A_{\mathrm{desc}}[1],
  \]
  where
  \[
    \bar A=\bigoplus_{R+S>0}\C\cdot[x^Ry^S],
    \qquad
    \bar A_{\mathrm{desc}}[1]
      =\bigoplus_{a+b>0}\C\cdot[\Psi_{a,b}].
  \]
  The Hamiltonian comparison uses only
  \(\bar A\oplus\bar A_{\mathrm{desc}}[1]\), and recognises this latter
  pair as the special-fibre, \(\hbar=0\), shadow of the Rees PBW lattice
  of Construction~\ref{constr:app-full-pbw-special-fibre-basis}: the
  closed Hamiltonian summand is the polynomial part
  \(\mathbf S(\mathfrak h)|_{\hbar=0}\), and the descendant copy is the
  primitive one-\(\psi\) PBW label sector
  \(\mathbf S^1(\mathfrak h^\vee[1])\otimes\mathbf S(\mathfrak h)|_{\hbar=0}\).
  The complementary sectors are exactly the scalar unit \([1]\), the
  scalar odd primitive
  \(\Psi_{0,0}\), and the vanishing higher-\(\psi\) primitive groups
  \(H_p(K_{R,S})=0\) for \(p\geq2\).
\end{thm}

\begin{proof}
  The constant Hamiltonian has zero Hamiltonian vector field, so the
  reduced Hamiltonian algebra is \(\bar A=\C[x,y]/\C\cdot1\) degree by
  degree.  This removes the degree-zero scalar class \([1]\).  The odd
  scalar \(\Psi_{0,0}\) is the one-\(\psi\) descendant of the same removed
  constant direction; it is a stable primitive trace class, but it has no
  partner in \(\bar A_{\mathrm{desc}}[1]\).
  Theorem~\ref{thm:app-full-primitive-psi-homology} proves that no
  primitive classes remain in \(\psi\)-degree at least two.  What remains
  is exactly the positive-degree Hamiltonian and one-\(\psi\) PBW label
  sector displayed above.  The identification with the special-fibre
  Rees lattice is
  Construction~\ref{constr:app-full-pbw-special-fibre-basis}: the closed
  Hamiltonian summand is the polynomial special fibre of the Rees algebra
  of \(\mathfrak h\), and the one-\(\psi\) descendant copy is the
  bidegree-graded \(\widetilde\Psi_{a,b}\) sector of the
  cotangent-shifted \(\mathfrak h^\vee[1]\) tensor.  This is a
  vector-space homology statement; it does not identify
  \(\Psi_{a,b}\) or its symmetrised normalisation
  \(\widetilde\Psi_{a,b}\) with Matlis principal parts as
  \(\mathfrak h\)-modules.
\end{proof}

\begin{cor}[Special-fibre defect polarization]
\label{cor:app-full-defect-polarization-slogan}
  The polynomial one-\(\psi\) PBW labels \(\Psi_{a,b}\) and the Matlis
  principal parts \(\rho_{a,b}\) are joined by the bidegree-only
  associated-graded morphism
  \(\operatorname{gr}\Phi\) of
  Proposition~\ref{prop:app-full-gr-Phi}: this morphism intertwines the
  two index sets at the special fibre \(\hbar=0\) and only there.  In
  particular, the full primitive \(\psi\)-homology of this appendix
  computes the special-fibre shadow of the cotangent module on the open
  brane sector, while the Matlis principal-parts module of
  \S~\ref{sec:pro-matlis-target} is the deformation-level
  cotangent module itself.  The formal-stalk chart measures positions
  polynomially through \(\bar A\), and its cotangent momenta live as
  residues through \(\mc P\); the shared bidegree label is
  associated-graded data, not module data.
\end{cor}

\begin{proof}
  Construction~\ref{constr:app-full-pbw-special-fibre-basis} identifies
  \(\widetilde\Psi_{a,b}\) with the unit-normalised \(\hbar=0\) class in
  \(\mathbf S^1(\mathfrak h^\vee[1])\otimes\mathbf S(\mathfrak h)|_{\hbar=0}\)
  in bidegree \((a,b)\).  The bidegree-only relabelling
  \(\operatorname{gr}\Phi\) of
  Proposition~\ref{prop:app-full-gr-Phi} is an isomorphism of bigraded
  vector spaces and not of \(\mathfrak h\)-modules; the obstruction is
  Theorem~\ref{thm:app-matlis-polynomial-descendant-obstruction}.
  The Matlis side is the continuous coadjoint module
  identified by Theorem~\ref{thm:app-matlis-residue-rigidity} as the
  unique perfect continuous coadjoint pairing.  Therefore
  \(\bar A_{\mathrm{desc}}[1]\) is the special-fibre shadow and
  \(\mc P\) is the deformation-level cotangent module, with the shared
  index set carrying associated-graded label data only.  This is the
  separation displayed by Theorem~\ref{thm:pbw-vs-deformation}.
\end{proof}
